\subsection{Use of Large Language Models (LLMs)}

In this work, we used LLM as a writing assistance tool to improve phrasing, grammar, and sentence clarity.

\subsection{Multi-step sampling result}

Table~\ref{tab:ablation_multi_step} presents the multi-step sampling results in terms of PESQ and ViSQOL for the Flow Matching stage in Flow2GAN. The results show that increasing the number of sampling steps improves generation quality.

\begin{table}[h!]
\centering
\caption{Multi-step results on LibriTTS dev set for the Flow Matching stage of Flow2GAN.}
\label{tab:ablation_multi_step}
\setlength{\tabcolsep}{6pt} % Default value: 6pt
% \renewcommand{\arraystretch}{1.2}
\resizebox{0.4\linewidth}{!}{
\begin{tabular}{l|cc}
\toprule
Inference steps & PESQ$\uparrow$ & ViSQOL$\uparrow$ \\
\midrule
1 & 2.279 & 2.222 \\
2 & 3.469 & 3.749 \\
4 & 3.788 & 4.156 \\
8 & 4.013 & 4.392 \\
16 & 4.166 & 4.553 \\
\bottomrule
\end{tabular}
}
\end{table}



\subsection{Model configuration}
\label{ssec:model_conf}

% Table~\ref{tab:} 

\begin{table}[h!]
\centering
\caption{Detailed model configurations of Flow2GAN with Mel-spectrogram and Encodec audio token conditioning.}
\label{tab:model_conf}
\setlength{\tabcolsep}{2pt} % Default value: 6pt
% \renewcommand{\arraystretch}{1.2}
\resizebox{1\linewidth}{!}{
\begin{tabular}{l|c|c}
\toprule
Configuration & Mel-spectrogram & Encodec audio token \\
\midrule
Input condition & N-FFT=1024,Hop=256,N-Mel=100 & Latent-dim=128,Hop=320\\
Time embedding dimension & 512 & 512 \\
Branch resolutions (N-FFT,Hop) & (512,256), (256,128), (128,64) & (640,320), (320,160), (160,80) \\
Branch embedding dimensions & 768,512,384 & 768,512,384 \\
Branch Layers & 8,8,8 & 8,8,8 \\
Feedforward hidden factor & 3 & 3 \\
Convolution kernel size & 7 & 7 \\
Condition encoder embedding dimension & 512 & 512 \\
Condition encoder layers & 4 & 4 \\
Loss scaling $\mathcal{S}(x)$, (N-FFT,Hop,N-filters) & (1024,256,256) & (1280,320,320) \\
\bottomrule
\end{tabular}
}
\end{table}

\subsection{Results comparison of Mel and log-Mel conditioning}

Table~\ref{tab:comp_mel_logmel} compares Flow2GAN results with Mel and log-Mel conditioning on the LibriTTS test set.

\begin{table}[h!]
\centering
\caption{Comparison of Mel and log-Mel conditioning on LibriTTS test set.}
\label{tab:comp_mel_logmel}
\setlength{\tabcolsep}{6pt} % Default value: 6pt
% \renewcommand{\arraystretch}{1.2}
\resizebox{0.85\linewidth}{!}{
\begin{tabular}{l|c|cccccc}
\toprule
Model & Condition & PESQ$\uparrow$ & ViSQOL$\uparrow$ & V/UV F1$\uparrow$ & Periodicity$\downarrow$ & FSD$\downarrow$ \\
\midrule
Flow2GAN, 1-step & Mel & \textbf{4.288} & 4.939 & \textbf{0.976} & \textbf{0.058} & \textbf{0.027} \\
Flow2GAN, 1-step & log-Mel & 4.189 & \textbf{4.957} & 0.975 & 0.063 & 0.028 \\
\midrule
 Flow2GAN, 2-step & Mel & \textbf{4.454} & 4.966 & 0.981 & 0.045 & \textbf{0.019} \\
  Flow2GAN, 2-step & log-Mel & 4.440 & \textbf{4.979} & \textbf{0.983} & \textbf{0.044} & 0.023 \\
\midrule
Flow2GAN, 4-step  & Mel & \textbf{4.484} & 4.966 & 0.984 & 0.039 & 0.017 \\
Flow2GAN, 4-step  & log-Mel & \textbf{4.484} & \textbf{4.986} & \textbf{0.985} & \textbf{0.037} & \textbf{0.016} \\
\bottomrule
\end{tabular}
}
\end{table}